% The template is adapted for CS229 project milestone
\documentclass{article}

\usepackage[preprint]{cs229}
\usepackage[utf8]{inputenc} % allow utf-8 input
\usepackage[T1]{fontenc}    % use 8-bit T1 fonts
\usepackage{hyperref}       % hyperlinks
\usepackage{url}            % simple URL typesetting
\usepackage{booktabs}       % professional-quality tables
\usepackage{amsfonts}       % blackboard math symbols
\usepackage{nicefrac}       % compact symbols for 1/2, etc.
\usepackage{microtype}      % microtypography
\usepackage{xcolor}         % colors
\usepackage{amsmath}
\usepackage{fancyhdr}       % for headers and footers
% Setup header
\pagestyle{fancy}
\fancyhf{} % clear all header and footer fields
\fancyhead[C]{CS229 Project Milestone}
\renewcommand{\headrulewidth}{0.4pt}

% Custom title format for one-page document without the project title
\renewcommand{\maketitle}{
  \begin{center}
    {\bf \large Testing Invariance and Generalization of Geometry-Aware Sharpness Minimization Methods}\\
    \vspace{0.1cm}
    {\bf Team Members:} Manav Dahra\\
    \vspace{0.1cm}
    {\bf Emails:} mav3rik@stanford.edu\\
    \vspace{0.2cm}
  \end{center}
}

\begin{document}

\maketitle

\section{Motivation}
% What problem are you tackling? Is this an application or a theoretical result?
Recent advances in deep learning optimization highlight the importance of sharpness-aware methods that seek flat minima to improve generalization. One such method, Sharpness-Aware Minimization (SAM) \citep{foret2021sam}, has shown promising results. However, it is sensitive to the choice of model parameterization, as its perturbation radius is defined in the Euclidean metric.

Adaptive SAM (ASAM) \citep{kwon2021asam} and Monge SAM (M-SAM) \citep{jacobsen2025msam} try to address these limitations by introducing scale invariance and geometry-aware metrics, respectively. However, existing evaluations of M-SAM are limited to a single fine-tuning task on ResNet-18 and a multi-modal alignment task; a systematic comparison of all three optimizers under controlled reparametrization across multiple architectures has not been done. 

Understanding how these methods behave under reparametrization is crucial for developing more robust optimization techniques, and this gap motivates our empirical work on benchmarking SAM, ASAM, and M-SAM on ResNet-18 and ViT-B/32 architectures with a focus on reparametrization invariance and generalization performance.

\section{Methods}

\subsection{Optimizers}
We implement three optimizers from scratch (SAM, ASAM, M-SAM) and compare them under controlled conditions.

\textbf{SGD} serves as the reference baseline: vanilla stochastic gradient descent with momentum and weight decay, trained with a cosine-annealing learning rate schedule.

\textbf{SAM} \citep{foret2021sam} augments SGD with a two-step update. In the first step the parameters are perturbed to the approximate local loss maximum within an $\ell_2$ ball of radius $\rho$:
\[
  \hat{\epsilon}(w) = \rho \cdot \frac{\nabla_w \mathcal{L}(w)}{\|\nabla_w \mathcal{L}(w)\|_2},
\]
and in the second step the base SGD update is applied at the perturbed point $w + \hat{\epsilon}(w)$.  Because the perturbation is defined in Euclidean space it is \emph{not} invariant to function-preserving rescalings of the weights.

\textbf{ASAM} \citep{kwon2021asam} addresses this by replacing the $\ell_2$ norm with an adaptive (parameter-scaled) norm.  The perturbation becomes
\[
  \hat{\epsilon}_\text{ASAM}(w) = \rho \cdot \frac{T_w^2 \nabla_w \mathcal{L}(w)}{\|T_w \nabla_w \mathcal{L}(w)\|_2},
\]
where $T_w = \mathrm{diag}(|w| + \eta)$ with a small floor $\eta=0.01$.  Multiplying the perturbation by $|w|^2$ makes it scale-invariant: if every weight is multiplied by a scalar $\alpha$, $\hat{\epsilon}_\text{ASAM}$ is unchanged.  Our implementation exactly follows this formula.

\textbf{M-SAM} \citep{jacobsen2025msam} takes a Riemannian geometry perspective.  The loss is viewed as a function on the model-function manifold; the Monge metric on that manifold yields a correction factor $1/\sqrt{1 + \|\nabla_w \mathcal{L}(w)\|_2^2}$ applied to the SAM perturbation:
\[
  \hat{\epsilon}_\text{MSAM}(w) = \frac{\rho}{\|\nabla_w \mathcal{L}(w)\|_2} \cdot \frac{1}{\sqrt{1 + \|\nabla_w \mathcal{L}(w)\|_2^2}} \cdot \nabla_w \mathcal{L}(w).
\]
This factor shrinks large-gradient perturbations, promoting geometry-aware and reparametrization-robust updates.

\subsection{Model and Dataset}
We use \textbf{ResNet-18} \citep{he2016resnet} adapted for CIFAR-10 \citep{krizhevsky2009cifar}: the standard ImageNet stem (7$\times$7 conv, stride 2, max-pool) is replaced with a 3$\times$3 stride-1 conv and an identity layer to retain spatial resolution for 32$\times$32 inputs.  The final fully-connected layer maps from 512 to 10 classes.  All experiments use a single random seed (seed 42) at this stage; multi-seed evaluation is planned for the final report.

\subsection{Training Protocol}
All models are trained for 200 epochs on CIFAR-10 with:
\begin{itemize}
  \item Batch size 256, learning rate 0.1, momentum 0.9, weight decay $5\times10^{-4}$;
  \item Cosine-annealing LR schedule (no warmup);
  \item Perturbation radius $\rho$ swept over $\{0.005, 0.01, 0.02, 0.03, 0.05\}$ for SAM and M-SAM; a single $\rho=0.5$ for ASAM (following the original paper's recommended scale).
\end{itemize}
For both SAM variants two forward–backward passes are performed per iteration.

\subsection{Evaluation Metrics}
We track final test accuracy, training accuracy, cross-entropy loss curves, and the \emph{divergence rate} $\Delta = \mathcal{L}_\text{test} - \mathcal{L}_\text{train}$ as a proxy for generalization gap.  For reparametrization invariance we measure the variance of final test accuracy across scaling factors $\alpha \in \{0.1, 0.5, 1.0, 2.0, 10.0\}$ applied to a function-preserving ReLU-layer rescaling.  Sharpness will be estimated via the Hutchinson stochastic trace estimator $\mathrm{tr}(H)/d$ with Rademacher vectors, which avoids forming the full Hessian.


\section{Preliminary Experiments and Results}

\subsection{Baseline Training (Experiment 1)}
We have completed the full baseline sweep (200 epochs each) for all three optimizer families on ResNet-18/CIFAR-10, producing 12 checkpoint files: SGD ($\rho{=}0$), SAM at five $\rho$ values, ASAM at $\rho{=}0.5$, and M-SAM at five $\rho$ values.

Table~\ref{tab:sgd_baseline} reports the SGD result and selected training-curve milestones.

\begin{table}[h]
\centering
\caption{SGD baseline on ResNet-18/CIFAR-10 (200 epochs, seed 42).}
\label{tab:sgd_baseline}
\begin{tabular}{lcccc}
\toprule
Optimizer & $\rho$ & Train Acc & Test Acc & Divergence Rate $\Delta$ \\
\midrule
SGD & 0.000 & 100.00\% & 95.45\% & 0.1773 \\
\bottomrule
\end{tabular}
\end{table}

The SGD model converges to near-zero training loss (0.0015) while achieving 95.45\% test accuracy, consistent with published results for ResNet-18 on CIFAR-10.  The divergence rate of 0.177 indicates moderate overfitting, which is expected given training to full convergence without early stopping.

Figure~\ref{fig:sgd_curve} summarises the learning dynamics: test accuracy rises steadily through the first 150 epochs, with the bulk of generalization gains occurring during the cosine-annealing decay phase (epochs 100–200), when test accuracy improves from $\sim$91\% to 95.5\%.

\begin{figure}[h]
\centering
\begin{minipage}{0.48\textwidth}
\centering
\textit{[SGD training/test accuracy curves — figure to be included in final report]}
\end{minipage}
\caption{SGD training and test accuracy over 200 epochs on ResNet-18/CIFAR-10.}
\label{fig:sgd_curve}
\end{figure}

\subsection{Infrastructure and Implementation Validation}
The training pipeline, data loading (HuggingFace CIFAR-10 cached locally), optimizer implementations, and checkpoint saving have all been validated.  The reparametrization utility (\texttt{apply\_relu\_reparam}) for ResNet-18 BasicBlocks has been implemented and tested: it scales BN1 affine parameters by $\alpha$ and compensates with $1/\alpha$ on the input channels of the subsequent convolution, exactly preserving the network function for any $\alpha > 0$.  The Hutchinson sharpness estimator and 1D loss landscape visualizer are also implemented and unit-tested.

ViT-B/32 experiments are currently paused: GELU is not exactly 1-homogeneous, so the reparametrization via simple weight scaling introduces a measurable output deviation ($\sim$0.06 per unit for $\alpha=2$). This issue is being investigated; a corrected formulation will be incorporated before the final report.


\section{Next Steps}

\begin{enumerate}
  \item \textbf{Comparative baseline analysis.}  Extract and aggregate final metrics (test accuracy, train loss, divergence rate) for all 12 completed checkpoints and produce a table comparing SAM, ASAM, and M-SAM against SGD across all $\rho$ values.  Identify the best $\rho$ for each optimizer.

  \item \textbf{Reparametrization invariance experiment (Experiment 2).}  Run the full sweep: for each of the four optimizers, apply the five scaling factors $\alpha \in \{0.1, 0.5, 1.0, 2.0, 10.0\}$ to the ResNet-18 initialization and train for 200 epochs.  Report per-optimizer variance of final test accuracy across $\alpha$ values and test the hypothesis that SAM is most sensitive, ASAM is partially invariant, and M-SAM is most invariant.

  \item \textbf{Flatness / sharpness analysis (Experiment 3).}  Load the trained checkpoints and compute: (a) the Hutchinson trace estimate $\mathrm{tr}(H)/d$ for each optimizer's best checkpoint; (b) 1D loss landscapes along random filter-normalized directions to visualize flatness.  Correlate sharpness estimates with generalization gap.

  \item \textbf{Multi-seed evaluation.}  Re-run the best configurations with seeds $\{0, 1, 2\}$ (in addition to seed 42) to obtain SEM estimates and statistically sound comparisons.

  \item \textbf{ViT experiments (stretch goal).}  Resolve the GELU reparametrization approximation and replicate the key experiments on ViT-B/32 to test whether geometry-aware invariance generalises across architectures.
\end{enumerate}


\section{Team Contributions}
\begin{itemize}
    \item \textbf{Manav Dahra}: End-to-end implementation of all four optimizers (SGD, SAM, ASAM, M-SAM); ResNet-18 and ViT-B/32 model adapters; CIFAR-10 data pipeline; training loop with two-step SAM closure support; reparametrization utility for ResNet-18 BasicBlocks; Hutchinson trace estimator and 1D loss landscape analysis; experiment runner scripts and YAML configuration system; baseline training sweep (12 checkpoint runs); initial results collection and this milestone writeup.
\end{itemize}

% References
\bibliographystyle{ACM-Reference-Format}
\bibliography{reference}
% If you don't have a .bib file, you can use the following manual format:
%\begin{thebibliography}{9}
%\bibitem{ref1}
%Author, A. et al. (Year). Title of paper. \textit{Conference/Journal}, pages.
%
%\bibitem{ref2}
%Author, B. et al. (Year). Title of paper. \textit{Conference/Journal}, pages.
%\end{thebibliography}

\end{document}